
% ------------------------------------------------
% Alban pendant les guerres de religion (1562-1598)
% Page originale : 45
% ------------------------------------------------

\begin{center}
    {\Large\itshape Alban pendant les guerres de religion (1562-1598)}
\end{center}

\vspace{1em}

Le pays ne semble guère en sûreté puisque ses maîtres ne se soucient pas de cette montagne en effervescence. En effet, Claude d'Auraison, évêque de Castres et seigneur d'Alban, vend en 1566 le « baillage d'Alban » pour la somme de 1300 livres à Jean Albert, marchand de Villefranche. Du coup, le susdit commerçant se déclare « seigneur baron d'Alban ». On se rappelle qu'en 1557, il s'était livré à une mascarade religieuse au château des Farguettes. Comment tout cela se concilie-t-il avec l'acte par lequel le roi de France a arrenté les revenus d'Alban à Antoine Aussaguel, habitant d'Albi qui s'est déclaré sur le champ « seigneur d'Alban » ? Serait-ce, qu'encore une fois, albigeois et villefranchois se partagent les fruits d'une baronnie dont ils sont les coseigneurs ?

\vspace{0.8em}

L'an 1565 est assez peu chargé de faits divers, du moins connus. Mentionnons toutefois que les consuls d'Alban ont fait arrêter deux habitants. Ceux-ci, préférant la justice d'église à celle du roi, en appellent aux tribunaux ecclésiastiques castrais. On ne sait pas le résultat de cette option. Le 28 avril de cette même année, un grand-père maternel d'Antoine Pujol de la Tapie était enterré à l'âge de cent vingt ans et enseveli dans l'église de Notre Dame d'Alban. Voici ce qu'un cahier manuscrit possédé par la famille Pujol de la Tapie racontait de l'événement... « Il fut à Notre Dame d'Alban le 26 avril 1565, s'y confessa et communia. C'était un dimanche. Le lundi matin 27 avril, venant de faire la prière au jardin, s'était assis au coin du feu sur un fauteuil où il ne fut pas plus tôt qu'il dit que quelque chose lui faisait mal et rendit l'âme dans le moment. Il fut enseveli dans l'église Notre Dame d'Alban près des fonts baptismaux le mardi vingt huit, en laquelle église il avait été baptisé, à ce qu'il disait, il y avait près
